\section{Introduction}\label{sec:intro}

Optimization models in process systems engineering, energy networks and
optimal control are often built from local pieces. A flowsheet links unit
models along streams, a multiperiod model links consecutive time stages, and a
network model links quantities at adjacent nodes. The objective is then a sum
of terms that each depend on a few variables, and the graph recording which
variables share a term often has small treewidth. For variables with finitely
many values, small treewidth is exactly what dynamic programming needs:
nonserial dynamic programming, bucket elimination and junction-tree message
passing find a global optimum in time exponential only in the width of a tree
decomposition \cite{BerteleBrioschi1972,Dechter1999,KollerFriedman2009}.

Mixed-integer nonlinear models have continuous variables and nonconvex
terms, and there the situation is less clear. Spatial branch-and-bound
solvers relax individual nonconvex terms by convex underestimators and branch
on variable domains \cite{TawarmalaniSahinidis2005,BelottiEtAl2009,
MisenerFloudas2014,BelottiEtAl2013}; they do not use a tree decomposition to
organize the search. Discretizing the continuous variables and running a
dynamic program gives neither a valid bound nor a running time that scales
well with the accuracy: a uniform grid of mesh $h$ on a unit interval has
$\Theta(1/h)$ points, so tables grow like $h^{-p}$ in the bag size $p$. Sparse
polynomial optimization exploits the same structure through linear and
semidefinite hierarchies \cite{WakiEtAl2006,Lasserre2006,BienstockMunoz2018};
the known bounds on their size are polynomial in the inverse accuracy
\cite{BienstockMunoz2018,KordaMagronRiosZertuche2025}.

Small treewidth alone does not make continuous nonconvex problems
tractable. Del Pia and Khajavirad
\cite[Theorems~1--3]{DelPiaKhajavirad2026} recently showed that
box-constrained nonconvex quadratic programming (box QP) is solvable in
strongly polynomial time on forests, but strongly NP-hard at treewidth two
even with small integer coefficients, and that quartic minimization is
strongly NP-hard on paths. This paper asks which additional quantitative information makes
tree decompositions useful for \emph{certified} global optimization of
mixed-integer problems, and what fails when it is absent.

\paragraph{Results.}\label{sec:intro-results}
We consider $\min\{F(x):x\in X\}$, with optimal value $\OPT$, where
$F=\sum_af_a(x_{S_a})$ is a sum of factors, $X$ is a box in which some
coordinates are integer, and a tree decomposition of the scopes $S_a$ with
maximum bag size $p$ is given. Two quantities govern the results: upper bounds
$L_i$ on the curvature of $F$ along each coordinate (for a quadratic,
$L_i=\max\{H_{ii},0\}$), and quadratic growth of $F$ away from a minimizer
$x^*$. \emph{Point growth} $F(x)-\OPT\ge g\norm{x-x^*}^2$ on $X$ gives the
condition number $\kappa=\max\{1,L/g\}$ with $L=\max_iL_i$. \emph{Weighted
growth} $F(x)-\OPT\ge\gamma\sum_iL_i(x_i-x_i^*)^2$ on $X$ gives
$\bar\kappa=\max\{1,1/\gamma\}$. Point growth implies weighted growth with
$\bar\kappa\le\kappa$, and $\bar\kappa$ does not change when continuous
coordinates are rescaled (Definition~\ref{def:growth}).

Choose finite grids in the coordinates and subtract, at each node $v$ of
coordinate $i$, the unary correction $L_iw^2/8$, where $w$ is the length of the
longest grid interval ending at $v$. The minimum of the corrected objective
over the grid equals the best vertex lower bound of Bajaj and Hasan
\cite{BajajHasan2020}, taken with per-coordinate curvature bounds, over all
cells of the product grid (Proposition~\ref{prop:cellwise}), so it is a lower
bound on $\OPT$. Because the corrections are unary, dynamic programming on
the given decomposition computes this minimum and all coordinate
min-marginals without enumerating cells. Grid intervals whose
endpoint min-marginals exceed the value of a known feasible point are
removed, and the grids of the successive stages form a \emph{path
certificate} (Theorem~\ref{thm:certificate}).

Growth enters only the running time. Under weighted growth, grids whose
spacing grows linearly away from the current corrected minimizer
(\emph{graded grids}), combined with this filtering, keep
$O(\theta^{-1}\log(n+2))$ nodes per coordinate at every stage, independently
of the accuracy, when the grading $\theta$ satisfies $8\bar\kappa\theta^2\le1$
(Lemma~\ref{lem:states}); filtered uniform grids can need order
$\sqrt{n\kappa}$ nodes per coordinate (Corollary~\ref{cor:uniformgrid}).
The algorithm CT (Algorithm~\ref{alg:ct}) tries
$\theta=2^{-\mu}$ for $\mu=2,3,\dots$ with a cap on the grid size, so the
growth constant need not be known. Below, $I$ denotes the input length and
$q\ge0$ an integer.

\begin{theorem}[Main results]\label{thm:intro-main}
Let $F$ be an explicit rational quadratic on a mixed box $X$, given as a sum
of factors with a tree decomposition of maximum bag size $p$, and let
$L_i=\max\{0,\partial_{ii}F\}$.
\begin{enumerate}[label=(\alph*)]
\item On every instance, CT$(2^{-q})$ returns a feasible rational point $\hat
x$ and a path certificate showing $\beta\le\OPT\le F(\hat x)\le\beta+2^{-q}$.
The certificate is checked by recomputing one dynamic program per stage in
exact arithmetic, and its validity uses no convexity, uniqueness or growth
(Theorems~\ref{thm:certificate} and~\ref{thm:approx}).
\item If $F$ has weighted growth with condition number $\bar\kappa$, every grid
on which CT$(2^{-q})$ builds tables has $O(\sqrt{\bar\kappa}\log(n+2))$ nodes per
coordinate, independently of $q$, and CT$(2^{-q})$ uses at most
$f(p,\bar\kappa)(I+q+1)^5$ bit operations, where
$f(p,t)=c_0(c_1p\sqrt t)^pt(1+\log_2t)^2$ for absolute
constants $c_0,c_1$. The certificate has size, and is checked in time, within
the same bound. Neither the growth constant nor $\bar\kappa$ is an input
(Lemma~\ref{lem:states} and Theorem~\ref{thm:approx}).
\item EX (Algorithm~\ref{alg:ex}), which runs CT$(2^{-q})$ for
$q=1,2,4,\dots$ and snaps each incumbent to a rational stationary point of a
nearby face of the box until an exact test accepts it, returns a global
minimizer and $\OPT$ on every instance, with a certificate
whose validity uses neither growth nor uniqueness. If $F$ has point growth
with condition number $\kappa$, EX uses at most $f_1(p,\kappa)(I+1)^{C_1}$
bit operations, for a computable function $f_1$ and an absolute constant
$C_1$ (Theorem~\ref{thm:exact}).
\end{enumerate}
\end{theorem}

In parameterized terms, part~(b) is a fixed-parameter bound in the
parameters $p$ and $\lceil\bar\kappa\rceil$, which CT neither receives nor
computes; here $\bar\kappa$ is the smallest weighted condition number of the
instance, which exists (Section~\ref{sec:setting} and
Remark~\ref{rem:fpt}). To the best of our knowledge, no earlier result gives
a running time of the form $f(p,\kappa)\poly(I+q)$ for certified approximation, or
$f(p,\kappa)\poly(I)$ for exact solution, of nonconvex or mixed-integer box
quadratic programs (Section~\ref{sec:related}). Under point growth, the
approximation result extends, with $\kappa$ in place of $\bar\kappa$, to
factors that are polynomials of fixed degree (Corollary~\ref{cor:poly}), but
continuous polynomial objectives can have irrational minimizers
(Example~\ref{ex:polylimits}).
Exactness itself needs no uniqueness. Every rational mixed box QP has
\emph{set growth}: $F(x)-\OPT\ge g_S\dist(x,\mathcal S)^2$ on $X$ for some
$g_S>0$, where $\mathcal S$ is the set of minimizers; we put
$\kappa_S=\max\{1,L/g_S\}$. A certified gap of $2^{-k}$ with
$k=\poly(I)+\max\{0,\log_2(1/g_S)\}$ lets the snapping step return an exact
minimizer (Theorem~\ref{thm:transfer} and Remark~\ref{rem:setgrowth});
point growth in part~(c) only bounds the time needed to reach this
accuracy. Table~\ref{tab:results} lists these results together with the
extensions described below.

\begin{table}[tb]
\centering
\small
\setlength{\tabcolsep}{2.5pt}
\caption{Main results (lower bounds: Theorem~\ref{thm:intro-lower}), with a
tree decomposition unless stated otherwise. Bounds are in bit operations,
except where they count table entries or maximum flows; ``exact'' marks the
bound for computing an exact minimizer. Certificates are valid without the
growth hypotheses, which enter only the running times. With
$\kappa(L',g')=\max\{1,L'/g'\}$, $\kappa_V=\kappa(L^V,g)$ in the row of
Theorem~\ref{thm:valuefn} and $\kappa_V=\kappa(L^+,g)$, with the reduced
curvatures $L^+$, in the row of Theorem~\ref{thm:cv}; in these two rows the
exact bounds with $\kappa_V$ assume a quadratic $F$ with point growth. The
function $f'$ is computable, $\kappa_{\mathcal K}$ is the condition number of
the value function on the core (Proposition~\ref{prop:cr-growth}), and the
degrees of the polynomials and $C_1$ are absolute constants, except in the
row of Theorem~\ref{thm:valuefn}, where the degrees depend on the time bounds
of the oracles. In the TU row,
$K=\max\{K_Z,s/\eta+1,r(5+2\sqrt{n_c\kappa_c})\}$, $n_c$ is the number of
continuous coordinates and $K_Z\le I$ bounds the number of values of a
discrete coordinate.}\label{tab:results}
\begin{tabular}{@{}>{\raggedright\arraybackslash}p{2.0cm}
>{\raggedright\arraybackslash}p{6.2cm}
>{\raggedright\arraybackslash}p{1.8cm}
>{\raggedright\arraybackslash}p{5.9cm}@{}}
\toprule
Result & Setting and hypothesis & Parameters & Bound\\
\midrule
Thm.~\ref{thm:approx} & mixed box QP; weighted growth &
$p,\bar\kappa$ & $f(p,\bar\kappa)(I+q+1)^5$\\
Thm.~\ref{thm:exact} & as above; point growth & $p,\kappa$ & exact
$f_1(p,\kappa)(I+1)^{C_1}$\\
Thm.~\ref{thm:valuefn} & value function, exact polynomial-time oracles;
point growth of $V$ & $p,\kappa_V$ &
$f(p,\kappa_V)\poly(\kappa_V,I+q)$; exact $f(p,\kappa_V)\poly(\kappa_V,I)$\\
Thm.~\ref{thm:cv} & private convex QP blocks, certified responses; point
growth of $V$ & $p,\kappa_V$ & as above; exact $2^{O(p)}\poly(I)$ if
$L^+=0$\\
Thms.~\ref{thm:cr-oracle}, \ref{thm:cr-search}, \ref{thm:cr-exact} &
residual concave in each coordinate, balanced; core of $k$ continuous
coordinates, no decomposition; core growth & $k,\kappa_{\mathcal K}$ & one maximum flow per
query; exact $8^k(1+\sqrt{k\kappa_{\mathcal K}})^k\poly(I)$\\
Thm.~\ref{thm:balanced} & balanced box QP, no decomposition; point growth &
$n,\kappa$ & $(n\kappa(I+q+1))^{O(1)}$; exact $(n\kappa(I+1))^{O(1)}$\\
Thm.~\ref{thm:tu-approx} & TU coupling, aligned data, exactly feasible
output; set growth, at most $r$ optimal values per coordinate &
$p,\kappa_c,s/\eta,r$ & $O(I+q)$ levels of $K^p$ table entries\\
Thm.~\ref{thm:cells} & finite optimal set; set growth & $p,r,\kappa_S$ &
$K_S^p\poly(I+q)$, $K_S=O(r\sqrt{n\kappa_S})$; exact
$K_S^p\poly(I+\log\kappa_S)$\\
Thm.~\ref{thm:endpointset}, Cor.~\ref{cor:facecsp} & $H_{ii}\le0$ for all
$i$, mixed box; no growth; output: a description of the optimal set & $p$ &
$2^{O(p)}\poly(I)$\\
Thm.~\ref{thm:diagdiscovery} & continuous box QP, diagonal Lagrangian
certificate; set growth; outputs also an exact description of the optimal
set & $p,\kappa_S$ & exact
$f'(p,\kappa_S)\poly(I)$\\
\bottomrule
\end{tabular}
\end{table}

\paragraph{Scope of the growth hypothesis.}
For a box QP, minimality alone confines nonconvexity: at every minimizer, the
Hessian block $H_{J_0J_0}$ of the set $J_0$ of continuous coordinates
strictly inside their bounds is positive semidefinite, so every direction of
negative curvature involves active bounds or integer coordinates. Growth adds
quantitative conditions (Lemma~\ref{lem:growthcert}). Locally, point growth
forces $H_{J_0J_0}\succeq2gI$, so $\kappa\ge2L/\lambda_{\min}(H_{J_0J_0})$ if
$J_0\ne\emptyset$, and weighted growth forces
$H_{J_0J_0}\succeq2\gamma\diag(L_i)_{i\in J_0}$. Globally, every other local
minimizer $x'$ satisfies $F(x')-\OPT\ge g\norm{x'-x^*}^2$, so
$\kappa\ge L\norm{x'-x^*}^2/(F(x')-\OPT)$. Within these conditions an
instance can have an indefinite Hessian, $\kappa\le\frac{21}4$ and
exponentially many strict local minima (Example~\ref{ex:family}). The
condition number can be exponential in $I$, already for $n=2$
(Section~\ref{sec:setting}). We do not claim that $\kappa$ or $\bar\kappa$
is moderate for the application models mentioned above; they enter only the
running-time bounds, never the validity of the output.

\paragraph{Lower bounds.}
Parts (a)--(c) of Theorem~\ref{thm:intro-lower} assume P${}\ne{}$NP, the
exponential-time hypothesis (ETH) that, for some $c>0$, no algorithm decides
3-SAT on $n'$ variables in time $2^{cn'}$ \cite{ImpagliazzoPaturiZane2001},
or its randomized version (rETH), which makes the same claim for algorithms
with error probability at most $1/3$ \cite{DellEtAl2014}; part (d) is
unconditional.

\begin{theorem}[Lower bounds]\label{thm:intro-lower}
\leavevmode
\begin{enumerate}[label=(\alph*)]
\item Unless P${}={}$NP, no algorithm computes, for every rational continuous
box QP with a supplied tree decomposition of maximum bag size three and a
unique minimizer, $\OPT$ within $2^{-q}$ in time polynomial in
$I+q+\log\kappa$. The hard instances have $\log\kappa=O(I)$
(Proposition~\ref{lim:prop:unique} and Corollary~\ref{lim:cor:nopolylog}).
\item Computing $\OPT$ within $\frac12$ is NP-hard for box QPs with
$\bar\kappa\le\kappa\le2$, continuous or all-binary, by a reduction from
maximum independent set; under ETH it cannot be done in time
$2^{o(p)}I^{O(1)}$ (Proposition~\ref{prop:lbwidth}).
\item Under rETH there are no computable $f$ and computable nondecreasing
unbounded $\psi\ge1$ such that some algorithm outputs, for every integer box QP
with a supplied tree decomposition of maximum bag size $p$ and a unique
minimizer, a feasible point
within $\frac12$ of $\OPT$ in time $f(p)(\kappa I)^{p/\psi(p)}$
(Proposition~\ref{prop:lbproduct}).
\item For $p\ge1$ and $\kappa\ge8p$, every deterministic algorithm that
accesses $F$ only through its values and returns a point with gap below
$Lp/\kappa$ for all functions on $[0,1]^p$ with upper coordinate curvature $L$, a
unique minimizer and point growth $L/\kappa$ makes at least
$(\kappa/(8p))^{p/2}-1$ evaluations on one of them
(Proposition~\ref{lim:prop:oracle}).
\end{enumerate}
\end{theorem}

Thus, unless P${}={}$NP, neither parameter can be dropped and the dependence
on $\kappa$ cannot be polylogarithmic; under ETH the dependence on $p$ is
exponential even for $\kappa\le2$; and under rETH the exponent of $\kappa$ in
a bound $f(p)\kappa^{a(p)}\poly(I)$ cannot be $o(p)$, already for integer box
quadratics, whereas Theorem~\ref{thm:intro-main}(b) has exponent $p/2+O(1)$.
Because $\bar\kappa\le\kappa$, the same holds for $\bar\kappa$. In the
value-oracle model of part~(d), CT attains the optimal exponent $p/2$ of
$\kappa$ (Remark~\ref{lim:rem:oracle}).

Section~\ref{sec:limits} proves further limits. The expanding-box chain of
length $m$ (Example~\ref{ex:chain}) has bag size three and $\kappa\le80$, and
CT certifies accuracy $2^{-q}$ on it in time polynomial in $m$ and $q$, but
the exact Bellman message for its last state has at least $2^{m-1}$ quadratic
pieces
(Proposition~\ref{lim:prop:messages}). Growth towards a continuum of
minimizers does not bound the grids produced by filtering, and under growth
towards an unknown optimal set, a deterministic algorithm that queries values
and derivatives needs at least of order $\varepsilon^{-n/2}$ queries for
accuracy $\varepsilon$, already when $\kappa_S=1$. A path of local affine
equalities makes the problem NP-hard at bag size three and $\kappa=1$, and
matching separator moments of any finite order gives no local error bound in
terms of negative curvature. As far as we know, the expanding-box chain, the
unique-minimizer hardness with explicit growth (a refinement of
\cite[Remark~2]{DelPiaKhajavirad2026}), the rETH bound (a quadratic encoding of multicolored clique
with a unique minimizer) and the moment obstruction are new; the other statements adapt
standard constructions, among them the running-sum encoding of Subset Sum of
Bienstock and Mu\~noz \cite[Appendix~A]{BienstockMunoz2018} (see also
\cite[Example~1.1]{CifuentesParrilo2016}).

\paragraph{Conditional recourse.}
A stiff convex term inflates the condition number even when the subproblem
it defines is easy. Exact minimization over some variables for fixed values
of the others (\emph{recourse}) preserves upper coordinate curvatures and
growth, so the certificate and the growth analysis apply to the value
function whenever it is a sum of factors that can be evaluated exactly in
polynomial time (Theorem~\ref{thm:valuefn}). Piecewise-affine optimal
responses of private convex quadratic blocks certify a reduction of the
coordinate curvatures of the value function that is the largest possible for
the block (Theorem~\ref{thm:cv} and Proposition~\ref{prop:vf-curv}), but
recourse does not reduce the parameter to the ratio of negative curvature to
growth in general (Proposition~\ref{prop:cv-limit}). For a quadratic whose recourse
problem is concave in each coordinate and submodular after sign changes, each
recourse value is one minimum cut with a flow certificate
(Theorems~\ref{thm:cr-oracle} and~\ref{thm:cr-search}), and for balanced box
quadratics the grid problem itself is a minimum cut, so no tree decomposition
is needed (Theorem~\ref{thm:balanced}). Parametric quadratic programming, cut
representations of submodular quadratic functions and threshold encodings of
ordered labels are classical; what is new is the curvature and growth
accounting that keeps them inside the certificate.

\paragraph{Coupling constraints.}
The certificate rests on rounding each coordinate independently, which breaks
under coupling constraints. Totally unimodular (TU) coupling constraints with
mesh-aligned data admit exactly feasible correlated rounding, which restores
the certificate and, for quadratics, exact output. Under set growth with
finitely many optimal values per coordinate, the algorithm is polynomial for
fixed bag size $p$ when the condition number $\kappa_c$ (built from a
curvature bound for the continuous variables and the set-growth constant),
the box widths measured in mesh units ($s/\eta$) and the number $r$ of optimal values per coordinate are
polynomially bounded (Theorem~\ref{thm:tu-approx}). Because its meshes are
uniform, it is not fixed-parameter tractable in $(p,\kappa_c)$
(Section~\ref{sec:tu-limits}), and its pseudopolynomial first grid cannot be
removed unless P${}={}$NP (Proposition~\ref{lim:prop:constraints}).

\paragraph{Several minimizers.}
One graded grid can need exponential time already with two isolated
minimizers (Proposition~\ref{prop:twocenters}). If the optimal set is finite,
unions of uniform cells keep $O(r\sqrt{n\kappa_S})$ nodes per coordinate,
where $r$ bounds the number of optimal values of a coordinate
(Theorem~\ref{thm:cells}). For quadratics with nonpositive Hessian
diagonal on a mixed box, a dynamic program over the box endpoints describes
the whole optimal set as a constraint problem on the same tree decomposition.
From it, uniqueness of the minimizer, the dimension of the optimal set and,
if it is finite, the number of minimizers are computed in
$2^{O(p)}\poly(I)$ bit operations
(Theorem~\ref{thm:endpointset} and Corollary~\ref{cor:facecsp}). This
extends Rosenberg's face description of the optimal set of a multilinear
function \cite{Rosenberg1972} and zero-residual descriptions of all optimal
labelings of discrete problems \cite{WainwrightJaakkolaWillsky2005,Werner2007}
to strictly concave coordinates and mixed boxes, in factored form. For
continuous box quadratics with a classical diagonal Lagrangian certificate of
optimality \cite{JeyakumarRubinovWu2006,LiWuQuan2015}, which is the same at
every minimizer, an exact minimizer and an exact description of the optimal
set are found in $f'(p,\kappa_S)\poly(I)$ bit operations, for a computable function
$f'$, without knowledge of $g_S$ (Theorem~\ref{thm:diagdiscovery}).

\paragraph{Organization.}
Section~\ref{sec:related} reviews related work, and Section~\ref{sec:setting}
fixes the model, the curvature bounds and the growth conditions.
Sections~\ref{sec:grids}--\ref{sec:exact} prove Theorem~\ref{thm:intro-main},
Sections~\ref{sec:recourse}--\ref{sec:optsets} treat the extensions of
Table~\ref{tab:results}, Section~\ref{sec:limits} proves
Theorem~\ref{thm:intro-lower} and the further limits,
Section~\ref{sec:computation} reports experiments, and
Section~\ref{sec:conclusion} lists open problems. The appendices follow the
order of the sections and contain the longer proofs and secondary results.
